%%%%%%%%%%%%%%%%%%%%%%%%%%%%%%%%%%%%%%%%%
% a0poster Portrait Poster
% LaTeX Template
% Version 1.0 (22/06/13)
%
% The a0poster class was created by:
% Gerlinde Kettl and Matthias Weiser (tex@kettl.de)
%
% This template has been downloaded from:
% http://www.LaTeXTemplates.com
%
% License:
% CC BY-NC-SA 3.0 (http://creativecommons.org/licenses/by-nc-sa/3.0/)
%
%%%%%%%%%%%%%%%%%%%%%%%%%%%%%%%%%%%%%%%%%

%----------------------------------------------------------------------------------------
%	PACKAGES AND OTHER DOCUMENT CONFIGURATIONS
%----------------------------------------------------------------------------------------

\documentclass[a0,portrait]{a0poster}\usepackage[]{graphicx}\usepackage[svgnames]{xcolor}
%% maxwidth is the original width if it is less than linewidth
%% otherwise use linewidth (to make sure the graphics do not exceed the margin)
\makeatletter
\def\maxwidth{ %
  \ifdim\Gin@nat@width>\linewidth
    \linewidth
  \else
    \Gin@nat@width
  \fi
}
\makeatother

\definecolor{fgcolor}{rgb}{0.345, 0.345, 0.345}
\newcommand{\hlnum}[1]{\textcolor[rgb]{0.686,0.059,0.569}{#1}}%
\newcommand{\hlstr}[1]{\textcolor[rgb]{0.192,0.494,0.8}{#1}}%
\newcommand{\hlcom}[1]{\textcolor[rgb]{0.678,0.584,0.686}{\textit{#1}}}%
\newcommand{\hlopt}[1]{\textcolor[rgb]{0,0,0}{#1}}%
\newcommand{\hlstd}[1]{\textcolor[rgb]{0.345,0.345,0.345}{#1}}%
\newcommand{\hlkwa}[1]{\textcolor[rgb]{0.161,0.373,0.58}{\textbf{#1}}}%
\newcommand{\hlkwb}[1]{\textcolor[rgb]{0.69,0.353,0.396}{#1}}%
\newcommand{\hlkwc}[1]{\textcolor[rgb]{0.333,0.667,0.333}{#1}}%
\newcommand{\hlkwd}[1]{\textcolor[rgb]{0.737,0.353,0.396}{\textbf{#1}}}%
\let\hlipl\hlkwb
\usepackage{tikz}
\usepackage{framed}
\makeatletter
\newenvironment{kframe}{%
 \def\at@end@of@kframe{}%
 \ifinner\ifhmode%
  \def\at@end@of@kframe{\end{minipage}}%
  \begin{minipage}{\columnwidth}%
 \fi\fi%
 \def\FrameCommand##1{\hskip\@totalleftmargin \hskip-\fboxsep
 \colorbox{shadecolor}{##1}\hskip-\fboxsep
     % There is no \\@totalrightmargin, so:
     \hskip-\linewidth \hskip-\@totalleftmargin \hskip\columnwidth}%
 \MakeFramed {\advance\hsize-\width
   \@totalleftmargin\z@ \linewidth\hsize
   \@setminipage}}%
 {\par\unskip\endMakeFramed%
 \at@end@of@kframe}
\makeatother

\definecolor{shadecolor}{rgb}{.97, .97, .97}
\definecolor{messagecolor}{rgb}{0, 0, 0}
\definecolor{warningcolor}{rgb}{1, 0, 1}
\definecolor{errorcolor}{rgb}{1, 0, 0}
\newenvironment{knitrout}{}{} % an empty environment to be redefined in TeX

\usepackage{alltt}
\usepackage{multicol} % This is so we can have multiple columns of text side-by-side
\columnsep=100pt % This is the amount of white space between the columns in the poster
\columnseprule=0pt % This is the thickness of the black line between the columns in the poster

%\usepackage[svgnames,table]{xcolor} % Specify colors by their 'svgnames', for a full list of all colors available see here: http://www.latextemplates.com/svgnames-colors


\usepackage{graphicx} % Required for including images
\graphicspath{{figures/}} % Location of the graphics files
\usepackage{booktabs} % Top and bottom rules for table
\usepackage[font=small,labelfont=bf]{caption} % Required for specifying captions to tables and figures
\usepackage{amsfonts, amsmath, amsthm, amssymb} % For math fonts, symbols and environments

\newtheorem{theorem}{Theorem}
\newtheorem{lemma}{Lemma}
\newtheorem{definition}{Definition}
\newtheorem{remark}{Remark}
\DeclareMathOperator*{\argmin}{arg\,min}

% \usepackage{fontspec}
% \defaultfontfeatures{Mapping=tex-text,Scale=MatchLowercase}
% \setmainfont{Times}

\usepackage{mathpazo} % math & rm
\linespread{1.1}        % Palatino needs more leading (space between lines)
\usepackage[defaultsans]{lato}
\usepackage{courier} % tt
\normalfont
\usepackage[T1]{fontenc}
\usepackage{palatino}
\usepackage{enumitem}
\usepackage{tcolorbox}
%\usepackage{natbib}

\usepackage{sectsty}
\allsectionsfont{\sffamily\LARGE}
\subsectionfont{\sffamily\large}
\captionsetup[figure]{labelfont={color=DarkSlateGray}}
\definecolor{MyBlue}{HTML}{033a57}
\definecolor{LBlue}{HTML}{0C9CE8}
\IfFileExists{upquote.sty}{\usepackage{upquote}}{}
\begin{document}

%----------------------------------------------------------------------------------------
%	POSTER HEADER
%----------------------------------------------------------------------------------------
\hspace{-5.6cm}\noindent\fcolorbox{MyBlue}{MyBlue}{%
\begin{minipage}[c]{0.80\linewidth}
\vspace{2cm}
\hspace{5cm}
\begin{minipage}[c]{0.77\linewidth}
\fontsize{2.8cm}{2.8cm}\selectfont \color{white} Estimating the long-term effect\\of
early treatment initiation
\\in Parkinson's disease
\\using observational data
\end{minipage}
\vspace{1cm}
\selectfont \color{white} 
\sffamily
\large
\hspace{5cm}\begin{tabbing}
\hspace{5cm}\fontseries{m}\selectfont Lieneke van den Heuvel, Luc J.W. Evers, Marjan J. Meinders, Bart Post,\\ \hspace{5cm}Anne M. Stiggelbout, Tom Heskes, Bastiaan R. Bloem, \textbf{Jesse H. Krijthe}
\end{tabbing}
\vspace{2cm}
\end{minipage}}
%
\hspace{5cm}\begin{minipage}[c]{0.25\linewidth}
\includegraphics[width=16cm]{logo-ru.pdf}

\vspace{1cm}

\hspace{0cm}\includegraphics[width=14cm]{radboudumc.pdf}

\vspace{1cm}

\includegraphics[width=16cm]{logo-lumc-black.pdf}
\end{minipage}

\vspace{2cm} % A bit of extra whitespace between the header and poster content

%----------------------------------------------------------------------------------------

\begin{multicols}{2} % This is how many columns your poster will be broken into, a portrait poster is generally split into 2 columns


\color{SaddleBrown} % SaddleBrown color for the introduction
\large

\color{DarkSlateGray}
\section*{Early treatment initiation in Parkinson's Disease}

Parkinson's disease (PD) is a neurodegenerative disorder, whose symptoms are treated using medication therapy in the large majority of patients. The long-term effects of the timing of initiation of PD medication are largely unknown. Both patients and physicians may choose to delay initiation of PD medication (such as Levodopa) for various reasons, including concerns regarding the toxicity of therapy and early development of treatment-related side effects. Large observational PD cohort studies are increasingly available. We want to use available data from these studies to estimate the long-term effect of the duration of the medication therapy early in the disease, to complement the available evidence from RCTs, which are expensive to replicate and may be less representative of the population of interest.
%\tcbset{}
\vspace{1cm}
\begin{tcolorbox}[colback=yellow!25!white,colframe=yellow!25!white,boxsep=0.75cm]
\textbf{Two main goals:}
\begin{enumerate}
\item Find out what evidence longitudinal observational data can bring to bear on the question of early medication treatment initiation.
\item Show that causal inference methods that take time-varying confounding into account are essential when estimating such effects.
\end{enumerate}
\end{tcolorbox}

\color{DarkSlateGray} 
\section*{Study Design and Data}
The effect estimate of interest is the \textbf{effect of the duration of PD medication therapy within the first two years of follow-up, on outcomes at year 2, 3 and 4}. We use data from the Parkinson's Progression Markers Initiative, which gathered data on 423 de novo patients, with a follow-up of up to 8 years. For this analysis, we used:
\begin{itemize}
\item 302 patients who started and did not stop using PD medication;
\item Bi-annual visits during the first two years (5 visits);
\begin{itemize}
\item Covariates for adjustment: age, education, time of diagnosis, MDS-UPDRS ratings, MoCA, MSEADL;
\end{itemize}
\item Outcome measurements at year 2, 3 and 4;
\begin{itemize}
\item Main: MDS-UPDRS Part III subscore (movement);
\item Secondary: Schwab \& England Activities of Daily Living scores and ratings for side-effects.
\end{itemize}
\end{itemize}

\section*{Time-varying Confounding}
We want to estimate what the effect of early treatment initiation is. However, treatment initiation is influenced by the current disease state: patients who are doing worse are more likely to start treatment. At the same time, the decision whether to initiate treatment affects the disease state at the next time point. These effect occur not just at baseline, but at every subsequent time-point. This is known as time-varying confounding (treatment-confounder feedback) and we need to adjust for it to correctly estimate the effect of the intervention of changing the start of treatment.

\vspace{2cm}

\usetikzlibrary{positioning,shapes.geometric}

\hspace{4cm}\scalebox{1.4}{
\begin{tikzpicture}[%
    ->,
    shorten >=2pt,
    >=stealth,
    node distance=1cm,thick,
    noname/.style={%
      circle,
      thick,
      minimum width=2em,
      minimum height=2em,
      draw
    }
  ]
    \node[noname] (1)                                             {$L_0$};
    \node[noname] (2) [right=of 1]                                {$L_1$};
    \node[noname] (3) [right=of 2] {$L_2$};
    \node[noname] (4) [right=of 3] {$L_3$};
\node[noname] (5) [right=of 4] {$L_4$};
    \node[noname] (6) [below=of 2]                                 {$T_1$};
    \node[noname] (7) [below=of 3]                                {$T_2$};
    \node[noname] (8) [below=of 4]              {$T_3$};
\node[noname] (9) [below=of 5]              {$T_4$};

    \path (1) edge                   node {} (2)
          (2) edge                   node {} (3)
          (3) edge                   node {} (4)
          (4) edge                   node {} (5)
          (6) edge                   node {} (7)
(7) edge                   node {} (8)
(8) edge                   node {} (9)
          (1) edge                   node {} (6)
(2) edge                   node {} (7)
(3) edge                   node {} (8)
(4) edge                   node {} (9)
          (6) edge                   node {} (2)
(7) edge                   node {} (3)
(8) edge                   node {} (4)
(9) edge                   node {} (5)
(6) edge [looseness=2.1, out=-30, in=-20] node {} (5)
(7) edge [looseness=1.9, out=-40, in=-30] node {} (5)
(8) edge [looseness=1.8, out=-50, in=-40] node {} (5);
  \end{tikzpicture}
}

\normalsize

\section*{Methods}
\vspace{0.5cm}
\begin{tcolorbox}[colback=LBlue!15!white,colframe=LBlue!15!white,boxsep=0.25cm]
\color{DarkSlateGray}
\subsection*{Simple Adjustment}
\begin{minipage}[t]{0.6\textwidth}
For comparison we consider a model that does not adjust for confounding, and returns a linear model of treatment duration on outcomes. A slight extension of this model also adjusts for the outcome at baseline.
\end{minipage}
%second column
\begin{minipage}[t]{0.4\textwidth}
\large
\vspace{-0.5cm}$$y = \beta_0 + \gamma \cdot \frac{\textrm{sum}(\bar{t})}{2} + \beta_1 y_{0}$$
\end{minipage}
\end{tcolorbox}

\vspace{0.5cm}
\begin{tcolorbox}[colback=LBlue!15!white,colframe=LBlue!15!white,boxsep=0.25cm]
\color{DarkSlateGray}
\subsection*{Inverse Probability of Treatment Weighting}
\begin{minipage}[t]{0.6\textwidth}
To adjust for time-varying confounding we weight the observations by the inverse of the probability of the treatment choice at each visit (normalized by the prior probability of the treatment), modelled using logistic regression models at each time-point.
\end{minipage}
%second column
\begin{minipage}[t]{0.4\textwidth}
\large
$$w = \prod_{k=1}^4 \frac{f(t_k | t_{k-1})}{f(t_k | l_{k-1}, t_{k-1})}$$
\end{minipage}
\end{tcolorbox}

\vspace{0.5cm}
\begin{tcolorbox}[colback=LBlue!15!white,colframe=LBlue!15!white,boxsep=0.25cm]
\color{DarkSlateGray}
\subsection*{Parametric g-formula}
\begin{minipage}[t]{0.6\textwidth}
We simulate the disease progression of all patients under different treatment regimes, and summarize the resulting outcomes using a marginal structural model. The disease progression at each time-point is modelled using linear regression models, taking into account the measurement history as covariates.
\end{minipage}
%second column
\begin{minipage}[t]{0.4\textwidth}
\large
$$f(y|\bar{t}, \bar{l}) \prod_{k=0}^4 f(l_k | \bar{l}_{k-1}, \bar{t}_{k})$$
\end{minipage}
\end{tcolorbox}
\vspace{-2cm}

\section*{Results}
  \includegraphics[width=1.0\linewidth]{different-models.pdf}
\captionof{figure}{Effect of an additional year of PD medication therapy on outcomes after two years estimated using different methods. Intervals indicate $\pm 2$ times the standard deviation of the estimator, estimated using $1000$ bootstrap replicates. For Modified Schwab \& England, higher scores correspond to improved disease outcomes, while for MDS-UPDRS part I and III, lower scores correspond to more beneficial outcomes. Note that adjusting for more of the confounding effects shifts the estimates to more beneficial outcomes for all three measures, indicating the importance of using these types of methods when analyzing this data.}

\vspace{1cm}

\definecolor{shadecolor}{rgb}{0.969, 0.969, 0.969}\color{fgcolor}
\includegraphics[width=1.0\linewidth]{main-outcomes.pdf} 
\captionof{figure}{Effect of one year of PD medication respectively Levodopa treatment ($n=82$) during the first two years of follow-up. Effect on MDS-UPDRS part III subscore measured at year 2, 3 and 4, for inverse probability of treatment weighting and the parametric g-formula. Lower scores correspond to better disease outcomes. Intervals indicate $\pm 2$ times the standard deviation of the estimator, estimated using $1000$ bootstrap replicates.}

\vspace{1cm}
\begin{tcolorbox}[colback=yellow!25!white,colframe=yellow!25!white,boxsep=0.75cm]
\section*{Conclusions}
\begin{enumerate}
\item Adjusting for time-varying confounding is important when studying interventions that happen over time.
\item We find no evidence for worse outcomes and tentative evidence for beneficial effects of early treatment initiation for our main outcome of interest.
\item We find no evidence of worse, nor improved secondary outcomes, due to high uncertainty in the estimates.
\end{enumerate}
\end{tcolorbox}




%----------------------------------------------------------------------------------------
%	REFERENCES
%----------------------------------------------------------------------------------------
% \small
% \begin{thebibliography}{1}
% % journal article
% \bibitem{Loog2014}
% Loog, M. (2016). Contrastive pessimistic likelihood estimation for semi-supervised classification. \textit{IEEE Transactions on Pattern Analysis and Machine Intelligence} \textbf{38(3)}, 462--475.
% \bibitem{Krijthe2016}
% Krijthe, J.H., Loog, M. (2016). The pessimistic limits of margin-based losses in semi-supervised learning. \textit{arXiv} \textbf{1612.08875}.
% \bibitem{Krijthe2017}
% Krijthe, J. H., \& Loog, M. (2017). Projected Estimators for Robust Semi-supervised Classification.\textit{ Machine Learning}, \textbf{106(7)}, 993–1008.
% \bibitem{Sion1958}
% Sion, M. (1958). On general minimax theorems.\textit{ Pacific J. Math}, \textbf{8(1)}, 171–176.
% \end{thebibliography}

\end{multicols}
\end{document}
